\chapter{A First Look at Class Field Theory}

Throughout the whole course, we studied two main topics:
\begin{itemize}
    \item Galois extensions $L:K$ and Galois group $\operatorname{Gal}(L/K)$;
    \item The ideal class group $C_K$, describing factorization of ideals in $\mathcal O_K$.
\end{itemize}
Class field theory connects these two subjects in the special case when $\operatorname{Gal}(L/K)$ is abelian. Very roughly, it describes finite abelian
extensions of $K$ in terms of ideals of $\mathcal O_K$ together with congruence conditions on generators of principal ideals.

We first define the Artin map on the free abelian group generated by the unramified prime ideals. We then interpret this group concretely using fractional ideals and recover the
ordinary ideal class group as a quotient. This leads naturally to the Hilbert class field, and
finally to ray class groups, where congruence conditions account for ramification.

The general proofs of class field theory are beyond the scope of this course. Our goal is to
understand the statements, the elementary constructions that lead to them, and several
important examples. 

Throughout this chapter, we will freely use the analogues over a general number field $K$ of the
prime-factorization results proved earlier over $\mathbb Q$.

\section{Decomposition Groups and Frobenius Automorphisms}

Let $L:K$ be a finite Galois extension of number fields with rings of integers $\mathcal O_L$, $\mathcal O_K$. Write
\[
G=\operatorname{Gal}(L/K).
\]

Let $P \lhd \mathcal O_K$ be a nonzero prime ideal and let $\mathfrak P\lhd \mathcal O_L$
be a prime ideal lying above $P$, i.e.
\[
\mathfrak P\cap \mathcal O_K=P.
\]
The quotient rings
\[
k(P):=\mathcal O_K/P, \qquad k(\mathfrak P):=\mathcal O_L/\mathfrak P
\]
are finite fields. Put
\[
q:=|k(P)|=N(P).
\]
If
\[
f:=f(\mathfrak P/P)=[k(\mathfrak P):k(P)]
\]
is the residue degree, then
\[
|k(\mathfrak P)|=N(\mathfrak P)=q^f.
\]

\begin{definition}
The {\bf decomposition group} of $\mathfrak P$ over $P$ is
\[
D(\mathfrak P/P):=\{\sigma\in G\mid \sigma(\mathfrak P)=\mathfrak P\}.
\]
Equivalently, it is the stabilizer of $\mathfrak P$ under the natural action of $G$ on the
prime ideals of $\mathcal O_L$ lying above $P$.
\end{definition}

Recall from Section \ref{sec-Galoisprime} that $G$ acts transitively on the prime ideals
of $\mathcal O_L$ lying above $P$. Let
\[
\mathfrak P_1=\mathfrak P,\mathfrak P_2,\dots,\mathfrak P_g
\]
be the distinct primes above $P$. Then
\[
P\mathcal O_L=\mathfrak P_1^e\cdots\mathfrak P_g^e,
\qquad
N(\mathfrak P_i)=N(P)^f.
\]
Taking absolute norms gives
\[
N(P)^{[L:K]}=N(P)^{efg},
\]
so
\[
\boxed{[L:K]=efg.}
\]
By the orbit--stabilizer theorem,
\[
|D(\mathfrak P/P)|=\frac{|G|}{g}=\frac{[L:K]}{g}=ef.
\]

We now explain the three integers $e,f,g$ directly in terms of the Galois group. The number $g$ comes from the orbit of a prime, $e$ from its inertia group, and $f$ from the induced action on its residue field.

Every $\sigma\in D(\mathfrak P/P)$ preserves $\mathfrak P$, so it induces an automorphism
of the residue field:
\[
\overline\sigma:k(\mathfrak P)\longrightarrow k(\mathfrak P),
\qquad
\alpha+\mathfrak P\longmapsto \sigma(\alpha)+\mathfrak P.
\]
Since $\sigma$ fixes $K$, this induced automorphism fixes $k(P)$. Hence we obtain a group
homomorphism
\[
D(\mathfrak P/P)\longrightarrow \operatorname{Gal}(k(\mathfrak P)/k(P)).
\]

\begin{definition}
The kernel
\[
I(\mathfrak P/P):=\ker\left(D(\mathfrak P/P)
\longrightarrow \operatorname{Gal}(k(\mathfrak P)/k(P))\right)
\]
is called the {\bf inertia group} of $\mathfrak P$ over $P$.
\end{definition}

We quote the following standard theorem on prime ideals in a finite Galois extension.

\begin{theorem}\label{thm-efg}
Let $L:K$ be a finite Galois extension and let $\mathfrak P$ lie above $P$. Let $e$, $f$, $g$ be as given above. Then the following sequence
\[
1\longrightarrow I(\mathfrak P/P)
\longrightarrow D(\mathfrak P/P)
\longrightarrow \operatorname{Gal}(k(\mathfrak P)/k(P))
\longrightarrow 1
\]
is exact. Consequently,
\[
|I(\mathfrak P/P)|=e,
\]
because $|D(\mathfrak P/P)|=ef$ and
$|\operatorname{Gal}(k(\mathfrak P)/k(P))|=f$. 
\end{theorem}

The important case for us is when $P$ is unramified, i.e. $e=1$. Then
$I(\mathfrak P/P)=\{id\}$ and therefore
\[
D(\mathfrak P/P)\cong \operatorname{Gal}(k(\mathfrak P)/k(P)).
\]
Since $k(\mathfrak P)$ is a finite extension of the finite field $k(P)$,
its Galois group is cyclic of order $f$ and is generated by the finite-field Frobenius
\[
x\longmapsto x^q,
\qquad q=|k(P)|=N(P).
\]
This motivates the following definition.

\begin{definition}\label{def-Frob}
Assume that $P$ is unramified in $L$. The {\bf Frobenius automorphism} of
$\mathfrak P$ over $P$ is the unique element
\[
\operatorname{Frob}_{\mathfrak P}\in D(\mathfrak P/P)
\]
such that
\[
\operatorname{Frob}_{\mathfrak P}(\alpha)
\equiv \alpha^{N(P)}\pmod{\mathfrak P}
\]
for every $\alpha\in\mathcal O_L$.
\end{definition}

This convention is sometimes called the {\bf arithmetic Frobenius} convention.

\begin{proposition}\label{prop-FrobConjugate}
Let $P$ be unramified in $L$ and let $\sigma\in G$. Then
\[
\operatorname{Frob}_{\sigma(\mathfrak P)}
=
\sigma\operatorname{Frob}_{\mathfrak P}\sigma^{-1}.
\]
\end{proposition}

\begin{proof}
Let $\tau=\operatorname{Frob}_{\mathfrak P}$. For every $\alpha\in\mathcal O_L$,
\[
\tau(\alpha)\equiv \alpha^{N(P)}\pmod{\mathfrak P}.
\]
Applying $\sigma$ gives
\[
\sigma\tau(\alpha)
\equiv \sigma(\alpha)^{N(P)}\pmod{\sigma(\mathfrak P)}.
\]
Replacing $\alpha$ by $\sigma^{-1}(\alpha)$, we obtain
\[
\sigma\tau\sigma^{-1}(\alpha)
\equiv \alpha^{N(P)}\pmod{\sigma(\mathfrak P)}.
\]
By the uniqueness in Definition \ref{def-Frob},
$\sigma\tau\sigma^{-1}=\operatorname{Frob}_{\sigma(\mathfrak P)}$.
\end{proof}

Thus for a general Galois extension, a prime $P$ determines a conjugacy class of
Frobenius elements. If $L:K$ is abelian, conjugation does nothing, and the Frobenius
depends only on $P$.

\begin{definition}
Let $L:K$ be a finite abelian extension and let $P$ be unramified. The element
\[
\left(\frac{L/K}{P}\right):=\operatorname{Frob}_{P}
\in\operatorname{Gal}(L/K)
\]
is called the {\bf Artin symbol} of $P$.
\end{definition}

\begin{proposition}\label{prop-orderFrob}
Let $L:K$ be finite abelian and let $P$ be unramified. Then
\[
\operatorname{ord}\left(\frac{L/K}{P}\right)=f.
\]
In particular,
\[
P \text{ splits completely in }L
\quad\Longleftrightarrow\quad
\left(\frac{L/K}{P}\right)=id.
\]
\end{proposition}

\begin{proof}
Under the isomorphism
\[
D(\mathfrak P/P)\cong \operatorname{Gal}(k(\mathfrak P)/k(P)),
\]
the Frobenius corresponds to the generator
$x\mapsto x^{N(P)}$ of a cyclic group of order $f$. Hence its order is $f$.
Since $P$ is unramified, Theorem \ref{thm-efg} gives $[L:K]=fg$. Thus $P$ splits
completely if and only if $g=[L:K]$, if and only if $f=1$, if and only if its Frobenius
is the identity.
\end{proof}

\section{Examples of Frobenius Map} \label{sec-frobeg}

\begin{example}[Quadratic fields]\label{eg-quadraticFrob}
Let
\(K = \mathbb{Q}\) and \(L=\mathbb Q(\sqrt d)\)
where $d$ is squarefree, and let
\[
\operatorname{Gal}(L/\mathbb Q)=\{id,\sigma\},
\qquad \sigma(\sqrt d)=-\sqrt d.
\]
Recall that a rational prime $p$ is unramified if and only if $p$ does not divide the discriminant of $L$. 

By Theorem \ref{thm-factorizationOfRationalPrimes}, 
\begin{center}
$(p)$ splits in $\mathcal O_L$ \quad $\Leftrightarrow$ \quad 
$x^2-d$ has a root modulo $p$ \quad $\Leftrightarrow$ \quad
\(\left(\frac d p\right)=1.
\)
\end{center}
Therefore
\[
\boxed{\left(\frac{L/\mathbb Q}{p}\right)
=
\begin{cases}
 id, & \left(\frac d p\right)=1,\\
 \sigma, & \left(\frac d p\right)=-1.
\end{cases}}
\]
\end{example}


\begin{example}[Cyclotomic fields]\label{eg-cyclotomicFrob}
Let
\(L=\mathbb Q(\zeta_m),
\)
where $\zeta_m$ is a primitive $m$-th root of unity. We have already seen that
\[
\operatorname{Gal}(L/\mathbb Q)\cong (\mathbb Z/m\mathbb Z)^*,
\]
where $a\in(\mathbb Z/m\mathbb Z)^*$ corresponds to the automorphism
\(\sigma_a(\zeta_m)=\zeta_m^a.
\)


By Example \ref{eg-cycloram},  $p$ is unramified in $L$ iff $p \nmid m$. Consider
$\sigma_p\in\operatorname{Gal}(L/\mathbb Q)$ given by
\[
\sigma_p(\zeta_m)=\zeta_m^p.
\]
Since $\mathcal O_L=\mathbb Z[\zeta_m]$, every $\alpha\in\mathcal O_L$ has the form
$\alpha=f(\zeta_m)$ with $f(x)\in\mathbb Z[x]$. Modulo a prime $\mathfrak P$ above $p$
we are in characteristic $p$, and therefore
\[
\sigma_p(\alpha)=f(\zeta_m^p)
\equiv f(\zeta_m)^p
=\alpha^p
\pmod{\mathfrak P}.
\]
Since $N((p))=p$, this is exactly the defining Frobenius congruence. Hence
\[
\operatorname{Frob}_p(\zeta_m)=\zeta_m^p.
\]
and hence
\[
\boxed{\left(\frac{L/\mathbb Q}{p}\right)
= p\ (\operatorname{mod} m)}
\]
in $\operatorname{Gal}(\mathbb Q(\zeta_m)/\mathbb Q)
\cong(\mathbb Z/m\mathbb Z)^*$.

By Proposition \ref{prop-orderFrob}, the residue degree of a prime above $p$ is
\[
f=\operatorname{ord}_m(p),
\]
the multiplicative order of $p$ in $(\mathbb Z/m\mathbb Z)^*$. In particular,
\[
p\text{ splits completely in }\mathbb Q(\zeta_m)
\quad\Longleftrightarrow\quad p\equiv 1\pmod m.
\]
\end{example}

\begin{example}
Take $m=5$. Then
\[
\operatorname{Gal}(\mathbb Q(\zeta_5)/\mathbb Q)
\cong (\mathbb Z/5\mathbb Z)^*\cong\mathbb Z_4.
\]
For a rational prime $p\neq 5$:
\begin{itemize}
    \item if $p\equiv1\pmod 5$, then $f=1$ and $p$ splits completely;
    \item if $p\equiv4\pmod 5$, then $f=2$, so $(p)$ is a product of two prime ideals,
    each with residue degree $2$;
    \item if $p\equiv2$ or $3\pmod 5$, then $f=4$, so $(p)$ remains prime.
\end{itemize}
Thus the congruence class of $p$ modulo $5$ determines the factorization of $(p)$ in
$\mathbb Z[\zeta_5]$.
\end{example}


We end with an example where $L:K$ is unramified for all prime ideals $P \lhd \mathcal O_K$. This will play an important role in Hilbert Class Field Theory:
\begin{example}\label{eg-FrobMinusFive}
Let $
K=\mathbb Q(\sqrt{-5}),\qquad L=K(i),
$
and write
$$
\operatorname{Gal}(L/K)=\{1,\sigma\},
\qquad \sigma(i)=-i.
$$

We first check that $L:K$ is unramified. Since $m_i(x) =
x^2+1,
$
whose discriminant is $-4$, every prime $P\nmid 2$ is unramified in
$L:K$. Thus the only prime that still needs to be considered is
$$
P_2=(2,1+\sqrt{-5}),
$$
the unique prime of $K$ above $2$.

To study $P_2$, observe that
$$
\sqrt5=i\sqrt{-5}\in L,
$$
so
$$
L=K(\sqrt5)
 =K\left(\frac{1+\sqrt5}{2}\right).
$$
Put
$$
\eta=\frac{1+\sqrt5}{2}.
$$
Then $\eta$ satisfies
$$
f(x)=x^2-x-1,
$$
whose discriminant is $5$. Since $P_2\nmid 5$, the prime $P_2$ is
unramified in $L:K$.

Moreover,
$$
\mathcal O_K/P_2\cong\mathbb F_2,
$$
and modulo $P_2$ we have
$$
f(x)\equiv x^2+x+1.
$$
This polynomial has no root in $\mathbb F_2$, so it is irreducible.
Hence $P_2$ is inert in the quadratic extension $L:K$, and hence all primes above $P$ are unramified.

Now we compute $\left(\frac{L/K}{P}\right)$ for all such $P$: For $P = P_2$, since it is inert in $L$ (and does not split), one has
$$
\left(\frac{L/K}{P_2}\right)=\sigma.
$$

Now let $P\nmid 2$, and let $\mathfrak P$ be a prime of $L$ above $P$.
By the defining property of the Frobenius element,
$$
\operatorname{Frob}_P(i)
\equiv i^{N(P)}
\pmod{\mathfrak P}.
$$
Since $P\nmid 2$, the elements $i$ and $-i$ are distinct modulo
$\mathfrak P$. Also $N(P)$ is odd, and therefore
$$
i^{N(P)}
=
\begin{cases}
i,&N(P)\equiv 1\pmod 4,\\
-i,&N(P)\equiv 3\pmod 4.
\end{cases}
$$
It follows that
$$
\boxed{
\left(\frac{L/K}{P}\right)=
\begin{cases}
\sigma,&P=P_2\text{ or }N(P)\equiv3\pmod4,\\[2mm]
1,&N(P)\equiv1\pmod4.
\end{cases}}
$$

Thus we have computed the Frobenius element for every prime of $K$.
\end{example}

\section{The Artin Map}

Let $L:K$ be a finite abelian extension, and let $S$ be the finite set of finite prime ideals
of $\mathcal O_K$ which ramify in $L$. For every $P\notin S$, the Artin symbol
\[
\left(\frac{L/K}{P}\right)=\operatorname{Frob}_P
\]
is defined.

At this stage, we do not yet need fractional ideals. Define $I_K(S)$ to be the free abelian
group generated by the prime ideals $P\notin S$. Thus an element of $I_K(S)$ is a formal
finite product
\[
\mathfrak a=\prod_{P\notin S}P^{n_P},
\qquad n_P\in\mathbb Z,
\]
with all but finitely many $n_P$ equal to zero. Multiplication is given by adding the
exponents. Because the factorization is unique, we may extend the Artin symbol
multiplicatively from prime ideals to all elements of $I_K(S)$.

\begin{definition}\label{def-ArtinMap}
The {\bf Artin map} for $L:K$ is the homomorphism
\[
\Theta_{L/K}:I_K(S)\longrightarrow\operatorname{Gal}(L/K)
\]
defined by
\[
\boxed{
\Theta_{L/K}\left(\prod_{P\notin S}P^{n_P}\right)
=
\prod_{P\notin S}
\left(\frac{L/K}{P}\right)^{n_P}.}
\]
\end{definition}

The product is finite, and it is unambiguous because $\operatorname{Gal}(L/K)$ is abelian.
Thus the Artin map packages all Frobenius elements into one homomorphism. 

A main problem of class field theory is to (1) understand the domain $I_K(S)$ of $\Theta_{L/K}$ in terms of ideal class group $C_K$; and (2) describe the kernel of $\Theta_{L/K}$.

\begin{example}\label{eg-ArtinQi}
Let
\[
L=\mathbb Q(i),
\qquad
K=\mathbb Q.
\]
Then
\[
\operatorname{Gal}(L/K)=\{1,\sigma\},
\qquad
\sigma(i)=-i.
\]
The only ramified rational prime is $2$, so $I_{\mathbb Q}(\{2\})$ is the free abelian group
generated by the odd rational prime ideals $(p) \lhd K$. We have seen in Section \ref{sec-frobeg} that
\[
\operatorname{Frob}_p=
\begin{cases}
1,&p\equiv1\pmod4,\\
\sigma,&p\equiv3\pmod4.
\end{cases}
\]
Hence, if $n$ is an odd positive integer with prime factorization
\(n=\prod_{p\mid n}p^{v_p(n)},
\) then
\[\boxed{
\Theta_{L/K}\left((n)\right) = \Theta_{L/K}\left(\prod_{p\mid n}(p)^{v_p(n)}\right)
=
\prod_{p\mid n}\operatorname{Frob}_p^{v_p(n)}
=
\begin{cases}
1,&n\equiv1\pmod4,\\
\sigma,&n\equiv3\pmod4.
\end{cases}}
\]
In particular,
\[
\Theta_{L/K}((3))=\sigma\neq1.
\]
\end{example}

Here is an example where $L:K$ is unramified:
\begin{example}
\label{eg-ArtinPrincipalMinusFive}
Let
$
K=\mathbb Q(\sqrt{-5})$, $L=\mathbb Q(i,\sqrt5)
$ with
$
\operatorname{Gal}(L/K)=\{1,\sigma\},
$ and
$$
\operatorname{Frob}_P=
\begin{cases}
\sigma,&P=P_2\text{ or }N(P)\equiv3\pmod4,\\
1,&\text{otherwise},
\end{cases}
$$
We claim that every nonzero principal integral ideal of $K$ maps to the
identity under the Artin map. Let
$$
\alpha=a+b\sqrt{-5}\in\mathcal O_K,
\qquad
\alpha\neq0,
$$
and put
$$
n=N((\alpha))
=|N_{K/\mathbb Q}(\alpha)|
=a^2+5b^2.
$$

Write
$$
(\alpha)=P_2^{e_2}\prod_{P\neq P_2}P^{e_P}.
$$
By the Frobenius computation above,
$$
\Theta_{L/K}((\alpha))=\sigma^E,
\qquad
E\equiv
e_2+
\sum_{\substack{P\neq P_2\\N(P)\equiv3\pmod4}}
e_P
\pmod2.
$$
We want to show that $E$ is even. To do so, write
$$
a=2^r a_0,\qquad b=2^r b_0,
$$
where $a_0,b_0$ are not both even. Since
$$
(\alpha)=(2)^r(a_0+b_0\sqrt{-5})
=P_2^{2r}(a_0+b_0\sqrt{-5}),
$$
the factor $P_2^{2r}$ does not affect $E$ modulo $2$. Thus we may assume
that $a$ and $b$ are not both even.

Suppose first that one of $a,b$ is even and the other is odd. Then
$$
n=a^2+5b^2\equiv1\pmod4.
$$
In particular $e_2=0$. Since
$$
n=\prod_{P\neq P_2}N(P)^{e_P},
$$
and every odd $N(P)$ is congruent to $1$ or $3$ modulo $4$, the
congruence $n\equiv1\pmod4$ implies that
$$
\sum_{\substack{P\neq P_2\\N(P)\equiv3\pmod4}}e_P
$$
is even. Hence $E$ is even.

Now suppose that $a$ and $b$ are both odd. Then
$$
n=a^2+5b^2\equiv6\pmod8.
$$
Thus
$$
e_2 =1, \qquad 
\frac n2\equiv3\pmod4.
$$
But
$$
\frac n2=\prod_{P\neq P_2}N(P)^{e_P},
$$
so the number of factors $N(P)\equiv3\pmod4$, counted with
multiplicity, is odd. Therefore
$$
E\equiv1+1\equiv0\pmod2.
$$

Thus in every case
$$
\boxed{\Theta_{L/K}((\alpha))=1.}
$$
i.e. every principal ideal lies in the kernel of the Artin map. Later we will see that this is precisely what should happen: $L$ is the
Hilbert class field of $K =\mathbb  Q(\sqrt{-5})$.
%The same is therefore true for every principal fractional ideal, since
%$$\Theta_{L/K}((\alpha)(\beta)^{-1}) =
%\Theta_{L/K}((\alpha)) \Theta_{L/K ((\beta))^{-1}=1.$$
%Hence

\end{example}




\section{Fractional Ideals and the Ideal Class Group}
To relate the Artin map $\Theta_{L/K}$ to the ideal class group $C_K$, we now give a concrete interpretation of
the group $I_K(S)$ when $S = \emptyset$, i.e. all ideals $\mathfrak P \lhd \mathcal O_L$ over a prime ideal $P \lhd \mathcal O_K$ are unramified. In such a case, we call $L:K$ an {\bf unramified} extension.

Recall that in Chapter 6, $C_K$ was defined using equivalence classes of nonzero integral ideals
$I,J\lhd \mathcal O_K$, where
\[
I\sim J
\]
if there exists nonzero $\alpha,\beta\in\mathcal O_K$ such that
\[
(\alpha)I=(\beta)J.
\]

We now introduce fractional ideals.

\begin{definition}
A {\bf fractional ideal} of $K$ is a nonzero $\mathcal O_K$-submodule
\[
\mathfrak a\subseteq K
\]
such that there exists a nonzero $d\in\mathcal O_K$ with
\[
d\mathfrak a\subseteq\mathcal O_K.
\]
\end{definition}

Thus every nonzero ideal $I \lhd \mathcal O_K$ is a fractional ideal. Conversely, if $\mathfrak a$ is a fractional ideal and
$d\mathfrak a\subseteq\mathcal O_K$, then
\[
I:=d\mathfrak a
\]
is an integral ideal and
\[
\mathfrak a=\frac1d I.
\]

For a nonzero ideal $J\subseteq\mathcal O_K$, define its {\bf inverse fractional ideal} by
\[
J^{-1}:=\{x\in K:xJ\subseteq\mathcal O_K\}.
\]
In particular, for a principal ideal $(d)$,
\[
(d)^{-1}=\frac1d\mathcal O_K.
\]
Thus, if $\mathfrak a=\frac1d I$, then
\[
\mathfrak a=(d)^{-1}I.
\]

Now factor the two integral ideals:
\[
I=\prod_P P^{a_P},
\qquad
(d)=\prod_P P^{b_P},
\]
with $a_P,b_P\geq0$. Hence
\[
\mathfrak a
=(d)^{-1}I
=\prod_P P^{a_P-b_P}.
\]
Therefore every nonzero fractional ideal has a unique expression
\[
\mathfrak a=\prod_P P^{n_P},
\qquad n_P\in\mathbb Z,
\]
where all but finitely many $n_P$ are zero.


\begin{definition}
Let 
$$I_K := \left\{\prod_P P^{n_P}\ \middle|\ n_P\in\mathbb Z,\ n_P=0\text{ for all but finitely many }P\right\}$$
denote the abelian group of all nonzero fractional ideals of $K$, with multiplication of ideals as the group operation.
\end{definition}
Recall the definition of $I_K(S)$. In the case when $L:K$ is unramified, one therefore has
$$I_K(S) = I_K(\emptyset) = I_K.$$

We next single out the principal fractional ideals.

\begin{definition}
Let
\[
P_K:=\{(\alpha)=\alpha\mathcal O_K\mid \alpha\in K^\times\}.
\]
Then $P_K$ is a subgroup of $I_K$, called the group of {\bf principal fractional ideals}.
\end{definition}

The quotient $I_K/P_K$ is exactly the ideal class group defined in Chapter 6.

\begin{proposition}\label{prop-classgroup-fractional}
There is a natural isomorphism
\[
\boxed{
C_K\cong I_K/P_K.
}
\]
\end{proposition}

\begin{proof}
Define
\[
\Phi:C_K\longrightarrow I_K/P_K
\]
by
\[
\Phi([I])=I P_K,
\]
where $I\lhd\mathcal O_K$ is a nonzero integral ideal.

We first check that $\Phi$ is well-defined. Suppose
\[
I\sim J.
\]
By definition, there exist nonzero $\alpha,\beta\in\mathcal O_K$ such that
\[
(\alpha)I=(\beta)J.
\]
Therefore, as fractional ideals,
\[
IJ^{-1}=\left(\frac{\beta}{\alpha}\right)
\]
is principal. Hence
\[
IP_K=JP_K.
\]

The map $\Phi$ is surjective. Indeed, let $\mathfrak a$ be any fractional ideal. Choose
$d\in\mathcal O_K\setminus\{0\}$ such that
\[
I:=d\mathfrak a\subseteq\mathcal O_K.
\]
Then $I$ is an integral ideal and
\[
\mathfrak a=\frac1d I.
\]
Since $(1/d)$ is principal,
\[
\mathfrak a P_K=IP_K.
\]

Finally, suppose
\[
\Phi([I])=\Phi([J]).
\]
Then
\[
IP_K=JP_K,
\]
so
\[
I=(\gamma)J
\]
for some $\gamma\in K^\times$. Write
\[
\gamma=\frac{\alpha}{\beta}
\]
for some nonzero $\alpha,\beta\in\mathcal O_K$. Then
\[
(\beta)I=(\alpha)J.
\]
Hence $I\sim J$, so
\[
[I]=[J].
\]
Thus $\Phi$ is injective.

Therefore
\(C_K\cong I_K/P_K.
\)
\end{proof}


%By unique factorization of fractional ideals, the abstract group $I_K(S)$ introduced in the previous section is now identified with the subgroup
%\[ I_K(S) = \left\{ \prod_P P^{n_P}\in I_K \;\middle|\; n_P=0\text{ for every }P\in S \right\}.\]
%In other words, $I_K(S)$ is precisely the group of fractional ideals whose prime factorization involves no prime in $S$. Thus the domain of the Artin map has acquired its usual concrete interpretation.

\section{The Hilbert Class Field}
We now consider the first case in class field theory - we construct a unramified field extension of $H_K : K$ (Hilbert class field), so that $I_K(S) = I_K$, and the kernel of the Artin map is precisely $P_K$. 

\begin{theorem}[Hilbert Class Field Theorem]\label{thm-HilbertClassField}
For every number field $K$, there exists a finite abelian Galois extension
\[
H_K:K,
\]
called the {\bf Hilbert class field} of $K$, with the following properties:
\begin{enumerate}
    \item every nonzero prime ideal $P\subseteq\mathcal O_K$ is unramified in $H_K$ (and hence $I_K(S) = I_K$);
    \item for each real embedding
\(\sigma:K\hookrightarrow\mathbb R,
\) (see Section \ref{sec-minkowski}), there are exactly $[H_K:K]$ embeddings
\(\widetilde{\sigma}:H_K\hookrightarrow\mathbb R
\) extending $\sigma$;
    \item every finite abelian extension of $K$ satisfying (1) and (2) is contained in $H_K$;
    \item the Artin map induces an isomorphism
    \[
    \boxed{C_K\ \cong\ \operatorname{Gal}(H_K/K)},
    \qquad
    [P]\longmapsto\operatorname{Frob}_P.
    \]
\end{enumerate}
\end{theorem}

We will not prove the Hilbert class field theorem. Its proof is one of the main results of
global class field theory. What matters here is the precise relation between the theorem and
the quotient $C_K=I_K/P_K$.

\begin{corollary}\label{cor-HilbertKernel}
For the Hilbert class field,
\[
\boxed{\ker\Theta_{H_K/K}=P_K.}
\]
Equivalently, a fractional ideal has trivial Artin symbol in $H_K$ if and only if it is
principal.
\end{corollary}

\begin{proof}
The Artin map
\[
\Theta_{H_K/K}:I_K\longrightarrow\operatorname{Gal}(H_K/K)
\]
induces, by Theorem \ref{thm-HilbertClassField}, an isomorphism
\[
\theta_{H_K/K}:I_K/P_K\xrightarrow{\sim}\operatorname{Gal}(H_K/K).
\]
Therefore its kernel is exactly $P_K$.
\end{proof}

\begin{corollary}\label{cor-HilbertDegree}
The degree of the Hilbert class field is the class number:
\[
\boxed{[H_K:K]=h_K.}
\]
\end{corollary}

\begin{proof}
By Theorem \ref{thm-HilbertClassField},
\[
\operatorname{Gal}(H_K/K)\cong C_K.
\]
Hence
\([H_K:K]=|\operatorname{Gal}(H_K/K)|=|C_K|=h_K.
\)
\end{proof}

Thus the class number has the interpretation
\[
\boxed{
h_K=\text{degree of the maximal everywhere-unramified abelian extension of }K.}
\]

\begin{corollary}\label{cor-principalSplit}
Let $P$ be a prime ideal of $\mathcal O_K$. Then
\[
\boxed{
P\text{ is principal}
\quad\Longleftrightarrow\quad
P\text{ splits completely in }H_K.}
\]
\end{corollary}

\begin{proof}
Under the Artin isomorphism,
\[
[P]\longmapsto\operatorname{Frob}_P.
\]
Now
\[
P\text{ is principal}
\quad\Longleftrightarrow\quad
[P]=e,
\]
while Proposition \ref{prop-orderFrob} gives
\[
P\text{ splits completely in }H_K
\quad\Longleftrightarrow\quad
\operatorname{Frob}_P=id.
\]
The two conditions are therefore equivalent.
\end{proof}

\begin{corollary}
The following are equivalent:
\begin{enumerate}
    \item $C_K=\{e\}$;
    \item $h_K=1$;
    \item $\mathcal O_K$ is a PID;
    \item $H_K=K$.
\end{enumerate}
\end{corollary}

\begin{example}\label{eg-HilbertMinusFive}
Let
\(K=\mathbb Q(\sqrt{-5})\),
\(L=\mathbb Q(i,\sqrt5).
\) Earlier we computed
\[
C_K\cong\mathbb Z_2,
\qquad
h_K=2.
\]
In Example \ref{eg-FrobMinusFive} we showed directly that $L:K$ is a quadratic extension
unramified at every finite prime of $K$. Since $K$ is imaginary quadratic, it has no
real places, so $L:K$ is unramified everywhere. Hence, by the Hilbert class field theorem,
\[
L\subseteq H_K.
\]
But \([L:K]=2=h_K=[H_K:K],
\) and therefore
\[
\boxed{H_K=L=\mathbb Q(i,\sqrt5).}
\]
Consequently, the Artin map induces an isomorphism
\[
\theta_{L/K}: I_K/P_K \cong C_K
\stackrel{\sim}{\longrightarrow}
\operatorname{Gal}(H_K/K).
\]
(In Example \ref{eg-ArtinPrincipalMinusFive}, we checked directly that principal ideals lie in $\ker\Theta_{L/K}$.) More precisely, the nonprincipal prime ideal
\([P_2]=[(2,1+\sqrt{-5})]
\)
represents the nontrivial class in $C_K$, and one has
$$\theta_{L/K}([P_2]) = \sigma.$$
\end{example}


The Fundamental Theorem of Galois Theory gives a further interpretation of the subgroup
structure of $C_K$. If
\[
K\subseteq M\subseteq H_K
\]
and $A\leq C_K$ corresponds to $M$ under
$C_K\cong\operatorname{Gal}(H_K/K)$, then
\[
\operatorname{Gal}(M/K)\cong C_K/A.
\]
Thus subgroups of the ideal class group describe the intermediate abelian extensions inside
the Hilbert class field.

\begin{remark}
Sections 9.1--9.5 form the unramified part of the story. The remaining sections explain how
congruence conditions enlarge this picture to arbitrary finite abelian extensions.
\end{remark}

\section{Ray Class Groups}

The Hilbert class field is governed by the quotient $I_K/P_K$. For a ramified abelian
extension, the example $\mathbb Q(i):\mathbb Q$ above shows that $P_K$ is too large to
be the kernel of the Artin map. The correction is to keep only those principal ideals which
admit generators satisfying specified congruence and sign conditions.

\begin{definition}
A {\bf modulus} of $K$ is a formal product
\[
\mathfrak m=\mathfrak m_0\mathfrak m_\infty,
\]
where
\[
\mathfrak m_0=\prod_P P^{n_P},
\qquad n_P\geq1,
\]
is a nonzero integral ideal with only finitely many primes occurring, and
$\mathfrak m_\infty$ is a product of distinct real places of $K$ (here real place means an embedding $K \hookrightarrow \mathbb{R}$).
\end{definition}

The finite part $\mathfrak m_0$ records finite primes where ramification is allowed. The
infinite part $\mathfrak m_\infty$ records real places which are allowed to become complex.

Let
\[
I_K(\mathfrak m)
=
\left\{\prod_P P^{n_P}\in I_K\;\middle|\; n_P=0
\text{ whenever }P\mid\mathfrak m_0\right\}.
\]
Thus $I_K(\mathfrak m)$ is the group of fractional ideals prime to the finite part
$\mathfrak m_0$.

\begin{definition}\label{def-rayCongruence}
Let $\alpha\in K^*$. We write
\[
\alpha\equiv1\pmod{\mathfrak m}
\]
if the following two conditions hold:
\begin{enumerate}
    \item for every $P^{n_P}\mid\mathfrak m_0$, one has
    \[
    v_P(\alpha-1)\geq n_P;
    \]
    \item for every real place $\sigma$ occurring in $\mathfrak m_\infty$, one has
    \[
    \sigma(\alpha)>0.
    \]
\end{enumerate}
Here $v_P$ denotes the exponent of $P$ in the prime-ideal factorization of a principal
fractional ideal. Equivalently, when $\alpha=a/b$ with $a,b\in\mathcal O_K$ both coprime to
$\mathfrak m_0$, the first condition says
\[
a\equiv b\pmod{P^{n_P}}
\]
for every $P^{n_P}\mid\mathfrak m_0$.
\end{definition}

Let $P_{K,\mathfrak m,1}$ be the subgroup of $I_K(\mathfrak m)$ consisting of principal
fractional ideals $(\alpha)$ with
\[
\alpha\equiv1\pmod{\mathfrak m}.
\]
Notice that the congruence condition forces $v_P(\alpha)=0$ for every
$P\mid\mathfrak m_0$, so $(\alpha)$ is indeed prime to $\mathfrak m_0$.

\begin{definition}
The {\bf ray class group modulo $\mathfrak m$} is
\[
\boxed{C_K(\mathfrak m):=I_K(\mathfrak m)/P_{K,\mathfrak m,1}.}
\]
\end{definition}

\begin{example}
Let $\mathfrak m=1$ be the trivial modulus, we have
\[
I_K(1)=I_K,
\qquad
P_{K,1,1}=P_K,
\]
and therefore
\[
\boxed{C_K(1)=C_K.}
\]
Thus the ordinary ideal class group is the first ray class group, with no congruence or sign
conditions imposed.
\end{example}

\begin{example}[Ray class groups of $\mathbb Q$]\label{eg-rayQ}
Let $K=\mathbb Q$ and take the modulus
\[
\mathfrak m=m\infty,
\]
where $\infty$ denotes the unique real place of $\mathbb Q$. Every fractional ideal of
$\mathbb Z$ is principal. If $\mathfrak a\in I_{\mathbb Q}(m\infty)$, then among the two
generators $\pm r$ of $\mathfrak a$ there is a unique positive generator $r>0$. Write
\[
r=\frac ab,
\qquad a,b\in\mathbb Z_{>0},
\qquad \gcd(ab,m)=1.
\]
Define
\[
\Phi:I_{\mathbb Q}(m\infty)\longrightarrow(\mathbb Z/m\mathbb Z)^*
\]
by
\[
\Phi(\mathfrak a)=ab^{-1}\pmod m.
\]
This is well-defined. Indeed, if $a/b=a'/b'$ and all four integers are coprime to $m$,
then $ab'=a'b$, so after reducing modulo $m$ and multiplying by $b^{-1}(b')^{-1}$ we get
\[
ab^{-1}\equiv a'(b')^{-1}\pmod m.
\]
It is clearly a surjective homomorphism.

Now $\mathfrak a=(a/b)$ lies in the kernel precisely when
\(a\equiv b\pmod m.
\)
Because $a/b>0$, this is exactly the condition
\[
\frac ab\equiv1\pmod{m\infty}.
\]
Hence
\[
\ker\Phi=P_{\mathbb Q,m\infty,1}.
\]
By the First Isomorphism Theorem,
\[
\boxed{C_{\mathbb Q}(m\infty)\cong(\mathbb Z/m\mathbb Z)^*.}
\]
\end{example}
Note that the infinite place is essential in the above example. Without the positivity condition, the same ideal has generators $r$ and $-r$, whose residue classes differ by $-1$ modulo $m$; one would instead
identify these two classes.

\section{Artin Reciprocity and the Existence Theorem}
By the example in the previous section,
\[
C_{\mathbb Q}(m\infty)\cong(\mathbb Z/m\mathbb Z)^* \cong \operatorname{Gal}(\mathbb Q(\zeta_m)/\mathbb Q)
.
\]
This is no coincidence -- indeed this is a special case of the following:
\begin{theorem}[Artin Reciprocity]\label{thm-ArtinReciprocity}
Let $L:K$ be a finite abelian extension. Then there exists a modulus $\mathfrak m$ containing
all finite and infinite places which ramify in $L$ such that the Artin map
\[
\Theta_{L/K}:I_K(\mathfrak m)\longrightarrow\operatorname{Gal}(L/K)
\]
is surjective and
\[
P_{K,\mathfrak m,1}\subseteq\ker\Theta_{L/K}.
\]
Consequently, $\Theta_{L/K}$ factors through the ray class group and gives a surjection
\[
\boxed{C_K(\mathfrak m)\twoheadrightarrow\operatorname{Gal}(L/K).}
\]
\end{theorem}

Among all moduli through which the Artin map factors, there is a unique minimal one under
divisibility. It is called the {\bf conductor} of $L:K$.

The converse to Artin reciprocity is the existence theorem. It is most useful in the
following subgroup form.

\begin{theorem}[Existence Theorem]\label{thm-ClassFieldExistence}
Let $\mathfrak m$ be a modulus of $K$, and let $H$ be a subgroup satisfying
\[
P_{K,\mathfrak m,1}\subseteq H\subseteq I_K(\mathfrak m).
\]
Then there exists a unique finite abelian extension $L_H:K$ inside a fixed algebraic closure,
with conductor dividing $\mathfrak m$, for which the Artin map induces an isomorphism
\[
\boxed{I_K(\mathfrak m)/H\ \cong\ \operatorname{Gal}(L_H/K).}
\]
In particular, if one takes
\[
H=P_{K,\mathfrak m,1}
\]
then $K_{\mathfrak m}:=L_{P_{K,\mathfrak m,1}}$
is called the {\bf ray class field}, and one has:
\[
C_K(\mathfrak m)\cong\operatorname{Gal}(K_{\mathfrak m}/K).
\]
\end{theorem}


\begin{example} 
For the trivial modulus $\mathfrak m=1$, the subgroup on the ideal side is
$P_{K,1,1}=P_K$. Hence
\[
K_1=H_K,
\qquad
C_K=C_K(1)
\cong\operatorname{Gal}(H_K/K).
\]
Thus the Hilbert class field theorem is the unramified special case of the general theory.
\end{example}

\begin{example}
Recall by Example \ref{eg-rayQ} that
\[
C_{\mathbb Q}(m\infty)\cong(\mathbb Z/m\mathbb Z)^* \cong \operatorname{Gal}(\mathbb Q(\zeta_m)/\mathbb Q)
\]
We study $C_{\mathbb Q}(m\infty) \to \operatorname{Gal}(\mathbb Q(\zeta_m)/\mathbb Q)$ directly -- recall that for $\mathfrak a=(a/b)\in I_{\mathbb Q}(m\infty)$, the Artin map corresponds to
\[
\mathfrak a\longmapsto ab^{-1}\pmod m.
\]
This is exactly the isomorphism constructed in Example \ref{eg-rayQ}, whose kernel is
$P_{\mathbb Q,m\infty,1}$. Therefore, we verified the isomorphism
\[
\boxed{K_{m\infty}=\mathbb Q(\zeta_m).}
\]
in Theorem \ref{thm-ClassFieldExistence} in this setup.
\end{example}

As a consequence, the class-field description over $\mathbb Q$ becomes the following
classical theorem.

\begin{theorem}[Kronecker--Weber Theorem]\label{thm-KroneckerWeber}
Every finite abelian extension $L:\mathbb Q$ is contained in a cyclotomic field:
\[
L\subseteq\mathbb Q(\zeta_m)
\]
for some $m\geq1$.
\end{theorem}
\begin{proof}
Let $L:\mathbb Q$ be a finite abelian extension. By Artin reciprocity
(Theorem \ref{thm-ArtinReciprocity}), there exists a modulus $\mathfrak m$
such that
\[
P_{\mathbb Q,\mathfrak m,1}
\subseteq
\ker\Theta_{L/\mathbb Q}.
\]
Since $\mathbb Q$ has only one real place, after enlarging the modulus (which makes $P_{\mathbb Q,\mathfrak m,1}$ smaller) if
necessary we may take
\[
\mathfrak m=m\infty
\]
for some positive integer $m$. Hence
\[
P_{\mathbb Q,m\infty,1}
\subseteq
H:=\ker\Theta_{L/\mathbb Q}.
\]
By the existence theorem (Theorem \ref{thm-ClassFieldExistence}), the subgroup
$H$ corresponds to the extension $L:\mathbb Q$, while
$P_{\mathbb Q,m\infty,1}$ corresponds to the full ray class field
$K_{m\infty}$. Since the correspondence between subgroups and class fields is
inclusion-reversing,
\[
P_{\mathbb Q,m\infty,1}\subseteq H
\qquad\Longrightarrow\qquad
L\subseteq K_{m\infty}.
\]
By the preceding example,
\(K_{m\infty}=\mathbb Q(\zeta_m).
\) Therefore
\[
L\subseteq\mathbb Q(\zeta_m)
\]
as stated in the theorem.
\end{proof}

\begin{theorem}[Quadratic Reciprocity]\label{thm-QuadraticReciprocity}
Let $p$ and $q$ be distinct odd rational primes. Then
\[
\boxed{
\left(\frac qp\right)
\left(\frac pq\right)
=
(-1)^{\frac{p-1}{2}\frac{q-1}{2}}.
}
\]
\end{theorem}

\begin{proof}
Put
\[
q^*=(-1)^{\frac{q-1}{2}}q,
\qquad
L=\mathbb Q(\sqrt{q^*}).
\]
For $p\neq q$, Example \ref{eg-quadraticFrob} gives
\[
\operatorname{Frob}_p(\sqrt{q^*})
=
\left(\frac{q^*}{p}\right)\sqrt{q^*}.
\tag{1}
\]

Now we apply Artin reciprocity. Note that
\[
\operatorname{disc}(L)=q^*,
\]
so the conductor of $L:\mathbb Q$ divides $q\infty$. Hence
Theorem \ref{thm-ArtinReciprocity} implies that
\[
P_{\mathbb Q,q\infty,1}
\subseteq
\ker\Theta_{L/\mathbb Q},
\]
where
\[
\Theta_{L/\mathbb Q}:
I_{\mathbb Q}(q\infty)
\longrightarrow
\operatorname{Gal}(L/\mathbb Q)
\]
is the Artin map. Therefore $\Theta_{L/\mathbb Q}$ factors through the quotient
\[
C_{\mathbb Q}(q\infty)
=
I_{\mathbb Q}(q\infty)/P_{\mathbb Q,q\infty,1}.
\]
That is, there is a surjective homomorphism
\[
\theta_{L/\mathbb Q}:
C_{\mathbb Q}(q\infty)\cong(\mathbb Z/q\mathbb Z)^*
\longrightarrow
\operatorname{Gal}(L/\mathbb Q) \cong \{\pm 1\}.
\]
Thus the action of $\operatorname{Frob}_p$ on $\sqrt{q^*}$ is given by a
nontrivial quadratic character of $(\mathbb Z/q\mathbb Z)^*$.

Since $(\mathbb Z/q\mathbb Z)^*$ is cyclic, it has a unique nontrivial
quadratic character, namely
\[
a\longmapsto\left(\frac aq\right).
\]
Therefore
\[
\operatorname{Frob}_p(\sqrt{q^*})
=
\left(\frac pq\right)\sqrt{q^*}.
\tag{2}
\]

Comparing (1) and (2),
\[
\left(\frac{q^*}{p}\right)=\left(\frac pq\right).
\]
Since
\[
q^*=(-1)^{\frac{q-1}{2}}q
\qquad\text{and}\qquad
\left(\frac{-1}{p}\right)=(-1)^{\frac{p-1}{2}},
\]
we obtain
\[
\left(\frac pq\right)
=
(-1)^{\frac{p-1}{2}\frac{q-1}{2}}
\left(\frac qp\right),
\]
which is quadratic reciprocity.
\end{proof}
